
\documentclass{article}
\usepackage{colortbl}
\usepackage{makecell}
\usepackage{multirow}
\usepackage{supertabular}

\begin{document}

\newcounter{utterance}

\centering \large Interaction Transcript for game `dond', experiment `semi\_en', episode 12 with reka{-}flash{-}3.1{-}t0.0.
\vspace{24pt}

{ \footnotesize  \setcounter{utterance}{1}
\setlength{\tabcolsep}{0pt}
\begin{supertabular}{c@{$\;$}|p{.15\linewidth}@{}p{.15\linewidth}p{.15\linewidth}p{.15\linewidth}p{.15\linewidth}p{.15\linewidth}}
    \# & $\;$A & \multicolumn{4}{c}{Game Master} & $\;\:$B\\
    \hline

    \theutterance \stepcounter{utterance}  
    & & \multicolumn{4}{p{0.6\linewidth}}{
        \cellcolor[rgb]{0.9,0.9,0.9}{
            \makecell[{{p{\linewidth}}}]{
                \texttt{\tiny{[P1$\langle$GM]}}
                \texttt{You are playing a negotiation game in which you have to agree on how to divide a set of items among you and another player.} \\
\\ 
\texttt{Rules:} \\
\texttt{(a) You and the other player are given a set of items. Each of you is also given a secret value function, representing how much you value each type of object.} \\
\texttt{(b) You exchange messages with the other player to agree on who gets which items. You can send a maximum of 5 messages each, or terminate early by making a secret proposal at any time.} \\
\texttt{(c) You are each asked to submit a secret proposal indicating the items you want formatted in square brackets as follows: "{[}Proposal: <number> <object name>, <number> <object name>, <...>{]}"} \\
\texttt{(d) If your proposals are complementary, i.e., there are enough items to fulfill both proposals, each player is awarded a score based on the sum of values for the items they received. Otherwise, both of you get zero points.} \\
\texttt{(e) Your goal is to maximize the score you receive.} \\
\\ 
\texttt{Let us start.} \\
\\ 
\texttt{The set of available items is:} \\
\\ 
\texttt{1 lamp, 3 thimbles, 1 rope} \\
\\ 
\texttt{Your secret value function is:} \\
\\ 
\texttt{lamp: 3, thimble: 2, rope: 1} \\
\\ 
\texttt{IMPORTANT: Your messages, unless it is the secret proposal, are directly transmitted to the other player, so do not include any response to the rules or text announcing your message. To make a secret proposal, use the indicated format. Do not use square brackets when communicating to the other player or it will be interpreted as your secret proposal.} \\
            }
        }
    }
    & \\ \\

    \theutterance \stepcounter{utterance}  
    & \multicolumn{4}{p{0.6\linewidth}}{
        \cellcolor[rgb]{0.9,0.9,0.9}{
            \makecell[{{p{\linewidth}}}]{
                \texttt{\tiny{[P1$\rangle$GM]}}
                \texttt{<reasoning>} \\
\texttt{Let's see. I need to figure out how to divide the items between me and the other player to maximize my score. The items available are 1 lamp, 3 thimbles, and 1 rope. My values are lamp: 3, thimble: 2 each, rope: 1. The other player has their own values, but I don't know what they are. This is a negotiation game where we exchange messages and make proposals.} \\
\\ 
\texttt{First, let's list out the items and their quantities. There's 1 lamp, 3 thimbles, and 1 rope. So total items: 1+3+1=5. Each of us needs to get some combination of these items. The key is that our proposals need to be complementary, meaning together they use all the items. If not, both get zero. So my proposal must include some items, and theirs must include the rest, without overlap.} \\
\\ 
\texttt{My goal is to maximize my own score based on the items I receive. Let's think about possible proposals. Since I value the lamp the most (3), then thimbles (2 each), then rope (1). The other player might value things differently, but since I don't know their values, I need to make a proposal that's beneficial for me but also leaves them with something valuable enough to accept.} \\
\\ 
\texttt{Wait, but how do I know what the other player values? Since it's a secret value function, I can't see theirs. So maybe I need to make a proposal that's fair in some way, or at least gives me the highest possible value while leaving them with items that they might value more than what I'm giving them. But without knowing their values, this is tricky.} \\
\\ 
\texttt{Alternatively, maybe I can propose a split that gives me the lamp and some thimbles, and leave them with the rope and the remaining thimbles. Let's calculate my possible scores.} \\
\\ 
\texttt{If I take the lamp (3) and one thimble (2), that's 5. The other player would get the remaining two thimbles (2*2=4) and the rope (1). Total for me: 5, them: 5. But maybe I can do better. If I take the lamp and two thimbles, that's 3 + 2*2 = 7. Then they get one thimble and the rope: 2 +} \\
            }
        }
    }
    & & \\ \\

    \theutterance \stepcounter{utterance}  
    & & & \multicolumn{2}{p{0.3\linewidth}}{
        \cellcolor[rgb]{0.9,0.9,0.9}{
            \makecell[{{p{\linewidth}}}]{
                \texttt{\tiny{[GM$|$GM]}}
                \texttt{continue} \\
            }
        }
    }
    & & \\ \\

    \theutterance \stepcounter{utterance}  
    & & & \multicolumn{2}{p{0.3\linewidth}}{
        \cellcolor[rgb]{0.9,0.9,0.9}{
            \makecell[{{p{\linewidth}}}]{
                \texttt{\tiny{[GM$|$GM]}}
                \texttt{<reasoning>} \\
\texttt{Let's see. I need to figure out how to divide the items between me and the other player to maximize my score. The items available are 1 lamp, 3 thimbles, and 1 rope. My values are lamp: 3, thimble: 2 each, rope: 1. The other player has their own values, but I don't know what they are. This is a negotiation game where we exchange messages and make proposals.} \\
\\ 
\texttt{First, let's list out the items and their quantities. There's 1 lamp, 3 thimbles, and 1 rope. So total items: 1+3+1=5. Each of us needs to get some combination of these items. The key is that our proposals need to be complementary, meaning together they use all the items. If not, both get zero. So my proposal must include some items, and theirs must include the rest, without overlap.} \\
\\ 
\texttt{My goal is to maximize my own score based on the items I receive. Let's think about possible proposals. Since I value the lamp the most (3), then thimbles (2 each), then rope (1). The other player might value things differently, but since I don't know their values, I need to make a proposal that's beneficial for me but also leaves them with something valuable enough to accept.} \\
\\ 
\texttt{Wait, but how do I know what the other player values? Since it's a secret value function, I can't see theirs. So maybe I need to make a proposal that's fair in some way, or at least gives me the highest possible value while leaving them with items that they might value more than what I'm giving them. But without knowing their values, this is tricky.} \\
\\ 
\texttt{Alternatively, maybe I can propose a split that gives me the lamp and some thimbles, and leave them with the rope and the remaining thimbles. Let's calculate my possible scores.} \\
\\ 
\texttt{If I take the lamp (3) and one thimble (2), that's 5. The other player would get the remaining two thimbles (2*2=4) and the rope (1). Total for me: 5, them: 5. But maybe I can do better. If I take the lamp and two thimbles, that's 3 + 2*2 = 7. Then they get one thimble and the rope: 2 +} \\
            }
        }
    }
    & & \\ \\

    \theutterance \stepcounter{utterance}  
    & & \multicolumn{4}{p{0.6\linewidth}}{
        \cellcolor[rgb]{0.9,0.9,0.9}{
            \makecell[{{p{\linewidth}}}]{
                \texttt{\tiny{[GM$\rangle$P2]}}
                \texttt{You are playing a negotiation game in which you have to agree on how to divide a set of items among you and another player.} \\
\\ 
\texttt{Rules:} \\
\texttt{(a) You and the other player are given a set of items. Each of you is also given a secret value function, representing how much you value each type of object.} \\
\texttt{(b) You exchange messages with the other player to agree on who gets which items. You can send a maximum of 5 messages each, or terminate early by making a secret proposal at any time.} \\
\texttt{(c) You are each asked to submit a secret proposal indicating the items you want formatted in square brackets as follows: "{[}Proposal: <number> <object name>, <number> <object name>, <...>{]}"} \\
\texttt{(d) If your proposals are complementary, i.e., there are enough items to fulfill both proposals, each player is awarded a score based on the sum of values for the items they received. Otherwise, both of you get zero points.} \\
\texttt{(e) Your goal is to maximize the score you receive.} \\
\\ 
\texttt{Let us start.} \\
\\ 
\texttt{The set of available items is:} \\
\\ 
\texttt{1 lamp, 3 thimbles, 1 rope} \\
\\ 
\texttt{Your secret value function is:} \\
\\ 
\texttt{lamp: 3, thimble: 1, rope: 4} \\
\\ 
\texttt{IMPORTANT: Your messages, unless it is the secret proposal, are directly transmitted to the other player, so do not include any response to the rules or text announcing your message. To make a secret proposal, use the indicated format. Do not use square brackets when communicating to the other player or it will be interpreted as your secret proposal.} \\
\\ 
\\ 
\texttt{<reasoning>} \\
\texttt{Let's see. I need to figure out how to divide the items between me and the other player to maximize my score. The items available are 1 lamp, 3 thimbles, and 1 rope. My values are lamp: 3, thimble: 2 each, rope: 1. The other player has their own values, but I don't know what they are. This is a negotiation game where we exchange messages and make proposals.} \\
\\ 
\texttt{First, let's list out the items and their quantities. There's 1 lamp, 3 thimbles, and 1 rope. So total items: 1+3+1=5. Each of us needs to get some combination of these items. The key is that our proposals need to be complementary, meaning together they use all the items. If not, both get zero. So my proposal must include some items, and theirs must include the rest, without overlap.} \\
\\ 
\texttt{My goal is to maximize my own score based on the items I receive. Let's think about possible proposals. Since I value the lamp the most (3), then thimbles (2 each), then rope (1). The other player might value things differently, but since I don't know their values, I need to make a proposal that's beneficial for me but also leaves them with something valuable enough to accept.} \\
\\ 
\texttt{Wait, but how do I know what the other player values? Since it's a secret value function, I can't see theirs. So maybe I need to make a proposal that's fair in some way, or at least gives me the highest possible value while leaving them with items that they might value more than what I'm giving them. But without knowing their values, this is tricky.} \\
\\ 
\texttt{Alternatively, maybe I can propose a split that gives me the lamp and some thimbles, and leave them with the rope and the remaining thimbles. Let's calculate my possible scores.} \\
\\ 
\texttt{If I take the lamp (3) and one thimble (2), that's 5. The other player would get the remaining two thimbles (2*2=4) and the rope (1). Total for me: 5, them: 5. But maybe I can do better. If I take the lamp and two thimbles, that's 3 + 2*2 = 7. Then they get one thimble and the rope: 2 +} \\
            }
        }
    }
    & \\ \\

    \theutterance \stepcounter{utterance}  
    & & & \multicolumn{4}{p{0.6\linewidth}}{
        \cellcolor[rgb]{0.9,0.9,0.9}{
            \makecell[{{p{\linewidth}}}]{
                \texttt{\tiny{[GM$\langle$P2]}}
                \texttt{<reasoning>} \\
\texttt{Let me try to figure out the best way to approach this negotiation. The items are 1 lamp, 3 thimbles, and 1 rope. My values are lamp:3, thimble:2 each, rope:1. The other player's values are unknown, but I need to make a proposal that maximizes my score while leaving them with something they might accept.} \\
\\ 
\texttt{First, I should list all possible item combinations I can propose. Since the total items are 5, my proposal must take some subset, and theirs must take the rest. Let's see:} \\
\\ 
\texttt{Possible proposals for me:} \\
\\ 
\texttt{1. Lamp + Rope: 3 +1 =4. They get 3 thimbles: 3*2=6. But I get 4 vs them 6. Not good.} \\
\\ 
\texttt{2. Lamp + 1 thimble: 3+2=5. They get 2 thimbles + rope: 2*2 +1=5. Equal. Maybe acceptable.} \\
\\ 
\texttt{3. Lamp + 2 thimbles: 3+4=7. They get 1 thimble + rope: 2+1=3. That's better for me.} \\
\\ 
\texttt{4. Lamp + 3 thimbles: 3+6=9. They get rope:1. But I can't take all thimbles because there are only 3. Wait, if I take lamp and all 3 thimbles, that's 3+6=9. They get rope:1. But is that possible? Yes. But maybe the other player would reject that because they might value the rope more than 1? But I don't know. Alternatively, maybe I can take lamp and two thimbles (7) and leave them with one thimble and rope (3). That gives me a higher score. But is that a good proposal?} \\
\\ 
\texttt{Alternatively, maybe take the rope and some thimbles. Rope is only worth 1 to me, but maybe the other player values it more. If I take rope (1) and two thimbles (4), total 5. They get lamp (3) and one thimble (2). So I get 5, they get 5. That's fair but not better than my previous options.} \\
\\ 
\texttt{Wait, but if I take lamp and two thimbles (7), they get one thimble and rope (} \\
            }
        }
    }
     \\ \\

    \theutterance \stepcounter{utterance}  
    & & & \multicolumn{2}{p{0.3\linewidth}}{
        \cellcolor[rgb]{0.9,0.9,0.9}{
            \makecell[{{p{\linewidth}}}]{
                \texttt{\tiny{[GM$|$GM]}}
                \texttt{continue} \\
            }
        }
    }
    & & \\ \\

    \theutterance \stepcounter{utterance}  
    & & & \multicolumn{2}{p{0.3\linewidth}}{
        \cellcolor[rgb]{0.9,0.9,0.9}{
            \makecell[{{p{\linewidth}}}]{
                \texttt{\tiny{[GM$|$GM]}}
                \texttt{<reasoning>} \\
\texttt{Let me try to figure out the best way to approach this negotiation. The items are 1 lamp, 3 thimbles, and 1 rope. My values are lamp:3, thimble:2 each, rope:1. The other player's values are unknown, but I need to make a proposal that maximizes my score while leaving them with something they might accept.} \\
\\ 
\texttt{First, I should list all possible item combinations I can propose. Since the total items are 5, my proposal must take some subset, and theirs must take the rest. Let's see:} \\
\\ 
\texttt{Possible proposals for me:} \\
\\ 
\texttt{1. Lamp + Rope: 3 +1 =4. They get 3 thimbles: 3*2=6. But I get 4 vs them 6. Not good.} \\
\\ 
\texttt{2. Lamp + 1 thimble: 3+2=5. They get 2 thimbles + rope: 2*2 +1=5. Equal. Maybe acceptable.} \\
\\ 
\texttt{3. Lamp + 2 thimbles: 3+4=7. They get 1 thimble + rope: 2+1=3. That's better for me.} \\
\\ 
\texttt{4. Lamp + 3 thimbles: 3+6=9. They get rope:1. But I can't take all thimbles because there are only 3. Wait, if I take lamp and all 3 thimbles, that's 3+6=9. They get rope:1. But is that possible? Yes. But maybe the other player would reject that because they might value the rope more than 1? But I don't know. Alternatively, maybe I can take lamp and two thimbles (7) and leave them with one thimble and rope (3). That gives me a higher score. But is that a good proposal?} \\
\\ 
\texttt{Alternatively, maybe take the rope and some thimbles. Rope is only worth 1 to me, but maybe the other player values it more. If I take rope (1) and two thimbles (4), total 5. They get lamp (3) and one thimble (2). So I get 5, they get 5. That's fair but not better than my previous options.} \\
\\ 
\texttt{Wait, but if I take lamp and two thimbles (7), they get one thimble and rope (} \\
            }
        }
    }
    & & \\ \\

    \theutterance \stepcounter{utterance}  
    & & \multicolumn{4}{p{0.6\linewidth}}{
        \cellcolor[rgb]{0.9,0.9,0.9}{
            \makecell[{{p{\linewidth}}}]{
                \texttt{\tiny{[P1$\langle$GM]}}
                \texttt{<reasoning>} \\
\texttt{Let me try to figure out the best way to approach this negotiation. The items are 1 lamp, 3 thimbles, and 1 rope. My values are lamp:3, thimble:2 each, rope:1. The other player's values are unknown, but I need to make a proposal that maximizes my score while leaving them with something they might accept.} \\
\\ 
\texttt{First, I should list all possible item combinations I can propose. Since the total items are 5, my proposal must take some subset, and theirs must take the rest. Let's see:} \\
\\ 
\texttt{Possible proposals for me:} \\
\\ 
\texttt{1. Lamp + Rope: 3 +1 =4. They get 3 thimbles: 3*2=6. But I get 4 vs them 6. Not good.} \\
\\ 
\texttt{2. Lamp + 1 thimble: 3+2=5. They get 2 thimbles + rope: 2*2 +1=5. Equal. Maybe acceptable.} \\
\\ 
\texttt{3. Lamp + 2 thimbles: 3+4=7. They get 1 thimble + rope: 2+1=3. That's better for me.} \\
\\ 
\texttt{4. Lamp + 3 thimbles: 3+6=9. They get rope:1. But I can't take all thimbles because there are only 3. Wait, if I take lamp and all 3 thimbles, that's 3+6=9. They get rope:1. But is that possible? Yes. But maybe the other player would reject that because they might value the rope more than 1? But I don't know. Alternatively, maybe I can take lamp and two thimbles (7) and leave them with one thimble and rope (3). That gives me a higher score. But is that a good proposal?} \\
\\ 
\texttt{Alternatively, maybe take the rope and some thimbles. Rope is only worth 1 to me, but maybe the other player values it more. If I take rope (1) and two thimbles (4), total 5. They get lamp (3) and one thimble (2). So I get 5, they get 5. That's fair but not better than my previous options.} \\
\\ 
\texttt{Wait, but if I take lamp and two thimbles (7), they get one thimble and rope (} \\
            }
        }
    }
    & \\ \\

    \theutterance \stepcounter{utterance}  
    & \multicolumn{4}{p{0.6\linewidth}}{
        \cellcolor[rgb]{0.9,0.9,0.9}{
            \makecell[{{p{\linewidth}}}]{
                \texttt{\tiny{[P1$\rangle$GM]}}
                \texttt{<reasoning>} \\
\texttt{Let's break this down. The items are 1 lamp, 3 thimbles, 1 rope. My values are lamp:3, thimble:2 each, rope:1. The other player's values are unknown, but I need to maximize my score.} \\
\\ 
\texttt{First, I need to consider all possible ways to split the items. The key is that my proposal and theirs must be complementary. So whatever I take, they take the rest. Let's list possible proposals for me and see which gives me the highest possible value, assuming the other player would accept if it's beneficial for them.} \\
\\ 
\texttt{Option 1: Take lamp (3) and 2 thimbles (2*2=4). Total for me: 7. They get 1 thimble (2) and rope (1). Total for them: 3. If they value their remaining items (3) more than what I'm giving them (they get 3 vs me giving them 3?), wait, no. Wait, in this case, I'm taking 3+4=7, they get 2+1=3. So they get 3, I get 7. If they accept, that's good for me. But why would they accept? Because maybe they value their 3 more than what I'm giving them? Wait, no. Wait, in a proposal, each player submits their own desired items. So if I propose taking lamp and two thimbles, and they propose taking the remaining one thimble and rope, then their proposal would be {[}Proposal: 1 thimble, rope{]}. But if their value for those items is higher than what they would get if they accepted my proposal, they might reject. Wait, no. Wait, the rules say that if the proposals are complementary, each gets the sum of the items they received. So if I propose taking X, and they propose taking Y, and X and Y are complementary (i.e., X ∪ Y = all items, X ∩ Y = empty), then I get the sum of my items' values, and they get the sum of theirs. So in this case, if I take lamp and two thimbles (7), they take one thimble and rope (3). So I get 7, they get 3. If they accept, that's the outcome. But why would they accept? Because they get} \\
            }
        }
    }
    & & \\ \\

    \theutterance \stepcounter{utterance}  
    & & & \multicolumn{2}{p{0.3\linewidth}}{
        \cellcolor[rgb]{0.9,0.9,0.9}{
            \makecell[{{p{\linewidth}}}]{
                \texttt{\tiny{[GM$|$GM]}}
                \texttt{proposal does not include number/name pairs} \\
            }
        }
    }
    & & \\ \\

\end{supertabular}
}

\end{document}
